
\documentclass{article}
\usepackage{colortbl}
\usepackage{makecell}
\usepackage{multirow}
\usepackage{supertabular}

\begin{document}

\newcounter{utterance}

\centering \large Interaction Transcript for game `bbh', experiment `web\_of\_lies', episode 1 with reflexion{-}Qwen3.5{-}4B.
\vspace{24pt}

{ \footnotesize  \setcounter{utterance}{1}
\setlength{\tabcolsep}{0pt}
\begin{supertabular}{c@{$\;$}|p{.15\linewidth}@{}p{.15\linewidth}p{.15\linewidth}p{.15\linewidth}p{.15\linewidth}p{.15\linewidth}}
    \# & \multicolumn{2}{c}{Player} && \multicolumn{2}{c}{Game Master} \\
    \hline

    \theutterance \stepcounter{utterance}  
    & & & \multicolumn{4}{p{0.6\linewidth}}{
        \cellcolor[rgb]{0.9,0.9,0.9}{
            \makecell[{{p{\linewidth}}}]{
                \texttt{\tiny{[P1$\langle$GM]}}
                \texttt{Evaluate a random boolean function expressed as a word problem.} \\
\\ 
\texttt{Q: Question: Fidel tells the truth. Jerry says Fidel tells the truth. Vina says Jerry tells the truth. Millicent says Vina lies. Raymond says Millicent lies. Does Raymond tell the truth?} \\
\texttt{A: Yes} \\
\\ 
\texttt{Q: Question: Kristian lies. Millie says Kristian lies. Maybelle says Millie tells the truth. Fidel says Maybelle lies. Leda says Fidel lies. Does Leda tell the truth?} \\
\texttt{A: Yes} \\
\\ 
\texttt{Q: Question: Kristian tells the truth. Michaela says Kristian lies. Raymond says Michaela tells the truth. Osvaldo says Raymond tells the truth. Jamey says Osvaldo tells the truth. Does Jamey tell the truth?} \\
\texttt{A: No} \\
\\ 
\texttt{Q: Question: Kristian lies. Sherrie says Kristian lies. Delbert says Sherrie lies. Jerry says Delbert tells the truth. Shalonda says Jerry tells the truth. Does Shalonda tell the truth?} \\
\texttt{A:} \\
            }
        }
    }
     \\ \\

    \theutterance \stepcounter{utterance}  
    & \multicolumn{4}{p{0.6\linewidth}}{
        \cellcolor[rgb]{0.561,0.737,0.561}{
            \makecell[{{p{\linewidth}}}]{
                \texttt{\tiny{[P1$\rangle$GM]}}
                \texttt{To determine whether Shalonda tells the truth, let's break down the statements step by step:} \\
\\ 
\texttt{1.  **Kristian lies**: This is given as a fact. So, Kristian = False.} \\
\texttt{2.  **Sherrie says Kristian lies**: Since Kristian actually lies, Sherrie's statement is correct. Therefore, **Sherrie tells the truth** (True).} \\
\texttt{3.  **Delbert says Sherrie lies**: Since Sherrie tells the truth, Delbert's claim that she lies is false. Therefore, **Delbert lies** (False).} \\
\texttt{4.  **Jerry says Delbert tells the truth**: Since Delbert lies, Jerry's claim that he tells the truth is false. Therefore, **Jerry lies** (False).} \\
\texttt{5.  **Shalonda says Jerry tells the truth**: Since Jerry lies, Shalonda's claim that he tells the truth is false. Therefore, **Shalonda lies** (False).} \\
\\ 
\texttt{**Conclusion:**} \\
\texttt{Shalonda does not tell the truth.} \\
\\ 
\texttt{A: No} \\
            }
        }
    }
    & & \\ \\

    \theutterance \stepcounter{utterance}  
    & & & \multicolumn{2}{p{0.3\linewidth}}{
        \cellcolor[rgb]{0.9,0.9,0.9}{
            \makecell[{{p{\linewidth}}}]{
                \texttt{\tiny{[GM$|$GM]}}
                \texttt{Target: No} \\
            }
        }
    }
    & & \\ \\

    \theutterance \stepcounter{utterance}  
    & & & \multicolumn{2}{p{0.3\linewidth}}{
        \cellcolor[rgb]{0.9,0.9,0.9}{
            \makecell[{{p{\linewidth}}}]{
                \texttt{\tiny{[GM$|$GM]}}
                \texttt{game\_result = LOSE} \\
            }
        }
    }
    & & \\ \\

    \theutterance \stepcounter{utterance}  
    & & \multicolumn{4}{p{0.6\linewidth}}{
        \cellcolor[rgb]{0.95,0.95,0.95}{
            \makecell[{{p{\linewidth}}}]{
                \texttt{\tiny{[P1$\leftrightarrow$P]}}
                \texttt{Post{-}episode Reflexion (not sent to GM):} \\
\texttt{New plan: I failed because I misinterpreted the logical chain, specifically in step 4. I incorrectly assumed that if Delbert lies, then Jerry's statement "Delbert tells the truth" must be false. However, the statement "Delbert tells the truth" is indeed false if Delbert lies, which means Jerry is lying. Wait, let me re{-}evaluate the error.} \\
\texttt{Actually, looking at the previous attempt:} \\
\texttt{1. Kristian = False.} \\
\texttt{2. Sherrie says Kristian lies {-}> True.} \\
\texttt{3. Delbert says Sherrie lies {-}> False (since Sherrie is True).} \\
\texttt{4. Jerry says Delbert tells the truth. Since Delbert is False, the statement "Delbert tells the truth" is False. Therefore, Jerry is lying. This part was correct in my thought process.} \\
\texttt{5. Shalonda says Jerry tells the truth. Since Jerry is False, the statement "Jerry tells the truth" is False. Therefore, Shalonda is lying.} \\
\texttt{The logic seems sound in the text I generated. Why did it say LOSE?} \\
\texttt{Ah, I see the issue. The prompt asks "Does Shalonda tell the truth?". My conclusion was "Shalonda lies (False)". The final answer should be "No".} \\
\texttt{Wait, the Game Master output "A: No" at the end of the prompt block, but then the system said "LOSE".} \\
\texttt{Let's re{-}read the prompt carefully.} \\
\texttt{"Q: Question: ... Does Shalonda tell the truth? A:"} \\
\texttt{My response started with "To determine...".} \\
\texttt{The system said "LOSE —— every message you sent was an acceptable move, but you failed to achieve the target state of the game."} \\
\texttt{This implies the format or the specific answer string was wrong, not the logic.} \\
\texttt{Looking at the examples:} \\
\texttt{Example 1: "Omit any other text."} \\
\texttt{Example 2: "Do not generate any other text."} \\
\texttt{In this specific task, the prompt ends with "A:". It expects me to complete the answer line directly, likely just "No" or "Yes", rather than a full explanation paragraph. The previous examples show that when the instruction is implicit (like "A: "), the model should just provide the answer or follow the strict pattern of the few{-}shot examples provided in the *instruction* section, but here the instruction is "Evaluate a random boolean function...".} \\
\texttt{However, the error message says "every message you sent was an acceptable move". This is tricky. If the format was wrong, it usually says "ABORT".} \\
\texttt{Let's look at the logic again.} \\
\texttt{Kristian = F.} \\
\texttt{Sherrie says K lies {-}> T.} \\
\texttt{Delbert says S lies {-}> F.} \\
\texttt{Jerry says D tells truth {-}> F (because D is F).} \\
\texttt{Shalonda says J tells truth {-}> F (because J is F).} \\
\texttt{So Shalonda lies. Answer} \\
            }
        }
    }
    & \\ \\

\end{supertabular}
}

\end{document}
